%% kapa: exam=1 round=19 year=2008 period=1 form=A subject=appraisal_law subject_ko=부동산관계법규 seq=25 number=65
%% source: Q-Net 감정평가사 기출문제 (cst00309) artlSeq=1001962 file=extracted/1008/1001962/2008년 제19회 감정평가사 1차시험 기출문제/1교시.hwp
%% conversion: kordoc 4.24.51 -> tools/gen1.py (auto); review=proofread-sonnet; flags=none
\begin{question}{65}
국유재산과 관련된 판례의 내용으로 옳지 않은 것은?
\begin{choices}
  \item 국․공유재산의 관리청이 행정재산을 사용·수익하는 자에 대하여 하는 사용․수익허가취소는 항고소송의 대상이 되는 행정처분이다.
  \item 국유잡종재산에 대한 시효취득을 부인하는 (구)국유재산법 제5조(국유재산의 보호) 제2항의 관련 규정은 헌법상의 평등원칙과 사유재산권 보장의 이념 및 과잉금지원칙에 반한다.
  \item 국유재산의 관리청이 행정재산을 사용·수익하는 자에 대하여 하는 사용료 부과는 사법상의 이행청구가 아니라 행정처분에 속한다.
  \item 국유재산의 무단점유로 인한 변상금징수권은 사법상의 권리채무를 내용으로 하는 것이다.
  \item 국립의료원 부설주차장 위탁관리용역운영계약은 강학상 특허에 해당한다.
\end{choices}
\end{question}
\answer{4}
